% titel: Blatt 3, A3: Eduktmengen für 5.00 g S2Cl2
% Teil II: der durchgerechnete Weg.

\subsection*{Gegeben und gesucht}
Ziel: $m(\mathrm{S_2Cl_2}) = \qty{5.00}{g}$. Gesucht sind die Massen an
Schwefel und Chlor, die dafür eingesetzt werden müssen. Ausgeglichene
Reaktion: $2\,\mathrm{S} + \mathrm{Cl_2} \rightarrow \mathrm{S_2Cl_2}$.

Molare Massen (gerundet): $M(\mathrm{S}) \approx \qty{32.07}{g/mol}$,
$M(\mathrm{Cl_2}) \approx \qty{70.90}{g/mol}$,
$M(\mathrm{S_2Cl_2}) = 2\cdot 32.07 + 2\cdot 35.45 \approx
\qty{135.04}{g/mol}$.

\subsection*{Ansatz}
Produktmenge in $n$ umrechnen; die benötigten Edukt-Stoffmengen folgen aus den
Stöchiometriefaktoren ($2$ für $\mathrm{S}$, $1$ für $\mathrm{Cl_2}$), dann
je in eine Masse umwandeln.

\subsection*{Rechnung}
\begin{align*}
  n(\mathrm{S_2Cl_2}) &= \frac{m}{M} = \frac{5.00}{135.04}
  \approx 3.70\times10^{-2}\ \mathrm{mol}, \\[2pt]
  n(\mathrm{S}) &= 2\,n(\mathrm{S_2Cl_2}) = 2\cdot 3.70\times10^{-2}
  = 7.40\times10^{-2}\ \mathrm{mol}, \\
  m(\mathrm{S}) &= 7.40\times10^{-2} \times 32.07
  \approx \qty{2.37}{g}, \\[2pt]
  n(\mathrm{Cl_2}) &= \phantom{2\,}n(\mathrm{S_2Cl_2})
  = 3.70\times10^{-2}\ \mathrm{mol}, \\
  m(\mathrm{Cl_2}) &= 3.70\times10^{-2} \times 70.90
  \approx \qty{2.63}{g}.
\end{align*}

\begin{ergebnis}
Für $\qty{5.00}{g}$ $\mathrm{S_2Cl_2}$ benötigt man
$m(\mathrm{S}) \approx \qty{2.37}{g}$ Schwefel und
$m(\mathrm{Cl_2}) \approx \qty{2.63}{g}$ Chlor.
\end{ergebnis}

\begin{kontrolle}
Masse der Edukte $2.37 + 2.63 = \qty{5.00}{g}$ entspricht exakt der
Produktmasse --- korrekt, da kein Stoff verloren geht. Querprobe über den
Schwefelanteil: $w(\mathrm{S}) = 2\cdot 32.07/135.04 = 0.475$, also
$0.475\cdot 5.00 = \qty{2.37}{g}$.
\end{kontrolle}

\begin{merksatz}
Eduktbedarf für eine Zielmenge: Produkt in Stoffmenge umrechnen und über die
Stöchiometriefaktoren auf die Edukte zurückrechnen --- nie Massen direkt
proportional setzen.
\end{merksatz}
